\documentclass[11pt]{article}
\usepackage[margin=1in]{geometry}
\usepackage{hyperref}
\usepackage{enumitem}
\usepackage{booktabs}
\usepackage{array}

\hypersetup{
    colorlinks=true,
    linkcolor=blue,
    urlcolor=blue,
    pdftitle={220119 - Lab 4 Assignment Links},
    pdfauthor={Ahsanul Haque (220119)}
}

\begin{document}

\begin{center}
    {\huge \textbf{AI \& ML Lab 4 -- Assignment Links}} \\[0.3cm]
    {\Large \textbf{Name: Ahsanul Haque}} \\[0.1cm]
    {\Large \textbf{ID: 220119}} \\[0.4cm]
    {\large \textbf{Decision Tree (CART vs ID3) \& K-Means Clustering}} \\[0.5cm]
\end{center}

\noindent All notebooks are configured to fetch their datasets from a public GitHub raw
URL, so the Colab links below can be opened and \textbf{Run All} will execute the full
pipeline with no manual file upload.

\section*{Repository and Notebooks}

\begin{enumerate}[leftmargin=*, itemsep=0.3cm]

    \item \textbf{GitHub Profile:} \\
    \url{https://github.com/ahsanjust/}

    \item \textbf{Lab 4 Folder in the Project Repository:} \\
    \url{https://github.com/ahsanjust/Artificial-Intelligence-and-Machine-Learning-Lab/tree/master/Lab_4}

    \item \textbf{Decision Tree (CART vs ID3) -- GitHub Notebook:} \\
    \url{https://github.com/ahsanjust/Artificial-Intelligence-and-Machine-Learning-Lab/blob/master/Lab_4/DT/220119_DT.ipynb}

    \item \textbf{Decision Tree (CART vs ID3) -- Google Colab:} \\
    \url{https://colab.research.google.com/drive/1DT_220119_Lab4_Assignment5}

    \item \textbf{K-Means Clustering -- GitHub Notebook:} \\
    \url{https://github.com/ahsanjust/Artificial-Intelligence-and-Machine-Learning-Lab/blob/master/Lab_4/KMeans/220119_KMeans.ipynb}

    \item \textbf{K-Means Clustering -- Google Colab:} \\
    \url{https://colab.research.google.com/drive/1KM_220119_Lab4_Assignment6}

\end{enumerate}

\section*{Performance Summary -- Decision Tree}

\begin{center}
\begin{tabular}{lccccc}
\toprule
\textbf{Model} & \textbf{Accuracy} & \textbf{Precision} & \textbf{Recall} & \textbf{F1-Score} & \textbf{AUC-ROC} \\
\midrule
CART (Gini)  & 1.0000 & 1.0000 & 1.0000 & 1.0000 & 1.0000 \\
ID3 (Entropy)& 1.0000 & 1.0000 & 1.0000 & 1.0000 & 1.0000 \\
\bottomrule
\end{tabular}
\end{center}

\noindent Both decision trees recover essentially the same splits on this dataset
because the target is highly separable by a small subset of features
(\texttt{anxiety\_level} and \texttt{stress\_level}).

\section*{K-Means Clustering Summary}

\begin{itemize}[leftmargin=*]
    \item \textbf{Dataset:} Teens Mental Health Dataset (1\,200 rows, 12 features
        after dropping the binary target).
    \item \textbf{Optimal K:} 4 (selected from the elbow in the WCSS curve with
        silhouette-based validation).
    \item \textbf{Cluster sizes:} 322, 267, 290, 321.
    \item \textbf{Real-world prediction:} 10 hand-crafted teen profiles were scaled
        with the same fitted \texttt{StandardScaler} used at training time and
        assigned to the 4 clusters by \texttt{model.predict}.
\end{itemize}

\section*{Cluster Interpretation}

\begin{center}
\begin{tabular}{cll}
\toprule
\textbf{Cluster} & \textbf{Profile} & \textbf{Size} \\
\midrule
0 & Older Male Teens with Elevated Addiction & 322 \\
1 & High-Stress Female Teens                & 267 \\
2 & Younger Male Teens with Low Anxiety     & 290 \\
3 & Balanced Female Teens with Low Stress   & 321 \\
\bottomrule
\end{tabular}
\end{center}

\vfill

\noindent {\small \textit{Note: All Colab notebooks have permissions set to ``Anyone
with the link can view.''}}

\end{document}
